% 01_introduction.tex —— 引言
\section{Introduction}

\subsection{Background}

China's recoverable tight-oil reserves are about 17.82 billion tonnes, roughly
40\% of its recoverable petroleum resources, and are concentrated in the Ordos,
Junggar, Sichuan and Songliao basins. Deep tight-oil matrices typically have
permeability below $0.1$~mD, small pore throats, and poorly connected natural
fractures; recovery under conventional development is commonly below $10\%$.
A composite fluid of liquid or supercritical CO$_2$ and a nano-imbibition agent
is used for stimulation and huff-n-puff development. The operation has three
stages: high-pressure injection and hydraulic fracturing, shut-in soaking, and
production. Soaking is the link between stimulation and economics: a short soak
leaves the CO$_2$ and the nano-agent confined near the fracture, while a long
soak defers production and accumulates operating and capital-occupancy costs.
Field and laboratory studies confirm that both the geomechanical response and
the nano--CO$_2$ synergy matter \cite{p1-03,p1-15}, and that live-oil and
dead-oil systems behave differently during soaking \cite{p2-02}.

\subsection{Problem Statement}

The problem asks for four connected results. Problem~1 requires the coupled
seepage--stress--damage equations, with assumptions, initial and boundary
conditions, and a clear distinction between the injection, soaking and
production stages. Problem~2 requires a dimensionless form, the controlling
parameter groups and their magnitudes, and an approximate analytical relation
between soak time and stimulation radius. Problem~3 is a bi-objective
optimization over the nano concentration, injection rate, injected volume,
CO$_2$ fraction, soak time and injection timing, subject to equipment, material,
pressure and duration constraints, for three reservoir scenarios and under
$\pm15\%$ parameter uncertainty. Problem~4 uses the supplied production data
for a data-driven model, evaluates net present value with the prescribed
economic parameters, compares the data-driven and mechanistic results, and
performs a global sensitivity analysis.

\subsection{Approach and Roadmap}

The governing model uses a one-dimensional domain normal to the fracture
and includes free and dissolved CO$_2$, nano adsorption, poroelasticity and
compressive-strain damage. Numerical tests use an insoluble three-phase
flow solver, split nano transport and separate local mechanics checks. Problem~2 derives the mirror-kernel pressure redistribution
from the injection-end profile and validates it against the three-phase
simulation. Problem~3 replaces the expensive forward model with a reduced
production model calibrated to the training groups of the supplied data, and
solves the bi-objective problem with NSGA-II. Problem~4 builds the grouped
data model, applies the Appendix~B cash-flow rules exactly at the 25~MPa and
300~K reference state, and reports Sobol sensitivities with the optimal soak
time as an output.

The same symbols and assumptions are used in all four problems. Where a
quantity is not identifiable from the supplied data---most importantly the
effect of injected volume and CO$_2$ fraction on production---we state the
assumption explicitly and report a structural sensitivity rather than
presenting it as a calibrated result.

\noindent\textit{AI provenance.} DeepSeek and OpenAI Codex assisted the
code, analysis and figures in P1--P4 and the sensitivity evaluation.
See the AI Tool Use Report for tool-specific contributions \cite{deepseek,codex}.
